\chapter{Сходимость случайных величин}

\section{Типы сходимостей случайных величин}

Пусть $ \xi,  \xi_1, \xi_2, \dotsc $ "--- случайные величины, заданные на одном вероятностном
пространстве.

\begin{definition}
    Будем говорить, что последовательность случайных величин $ \xi_1, \xi_2, \dotsc $
    \emph{сходится} к случайной величине $ \xi $ \emph{по вероятности} $ \xi_n \overarr{P} \xi $,
    если
    $$ \forall \eps > 0 \quad P\{|\xi_n - \xi| \ge \eps\} \underarr{n \to \infty} 0 $$
\end{definition}

\begin{definition}
    Последовательность случайных величин $ \xi_1, \xi_2, \dotsc $ \emph{сходится}
    к случайной величине $ \xi $ \emph{почти наверное} (\emph{с вероятностью 1})
    $ \xi_n \underarr{\text{п. н.}} \xi $, если
    $$ P\set{\lim \xi_n = \xi} = 1 $$

    \textcolor{gray}{Здесь имеется в виду поточечный предел функций $ \xi_n : \Omega \to \R $.}
\end{definition}

\begin{definition}
    Случайные величины $ \xi_1, \xi_2, \dotsc $ \emph{сходятся} к предельной случайной величине
    $ \xi $ \emph{в среднем порядка $ r $} $ \xi_n \overarr{r} \xi $, если существуют моменты
    $ \mathbb E \xi_n^r $ и $ \mathbb \xi^r $ порядка $ r $ и
    $$ \mathbb E |\xi_n - \xi|^r \underarr{n \to \infty} 0 $$
\end{definition}

\begin{definition}
    Будем говорить, что последовательность $ \xi_1, \xi_2, \dotsc $ \emph{сходится} к
    случайной величине $ \xi $ \emph{по распределению} $ \xi_n \overarr{d} \xi $, если
    для любой точки непрерывности $ x $ функции распределения $ F(x) $ предельной случайной величины
    $$ F_n(x) \underarr{n \to \infty} F(x) $$
\end{definition}

\begin{remark}
    \hfill
    \begin{itemize}
        \item \emph{почти наверное} $ \implies $ \emph{по вероятности}
        \item \emph{в среднем} $ \implies $ \emph{по вероятности}
        \item \emph{по вероятности} $ \implies $ \emph{по распределению}
    \end{itemize}
\end{remark}

\section{Центральная предельная теорема}

Начнём с самого простого варианта ЦПТ:
\begin{theorem}[Леви]
    $ \xi_1, \xi_2, \dotsc $ "--- последовательность одинаково распределённых случайных величин с
    математическим ожиданием $ a $ и ненулевой дисперсией $ \sigma^2 $.
    Тогда выполняется предельное соотношение
    $$ \sup\limits_x \Bigl| P \Set{\frac{\xi_1 + \dots + \xi_n - na}{\sigma \sqrt n} < x} -
    \Phi(x) \Bigr| \underarr{n \to \infty} 0 $$
\end{theorem}

\begin{proof}
    Отметим, что речь идёт об исходных случайных величинах, у которых существует момент второго
    порядка.
    НУО будем считать, что $ a = 0 $ и $ \sigma = 1 $.
    $$ S_n \coloneq \xi_1 + \dots + \xi_n $$

    Пусть $ f(t) $ "--- характеристическая функция величин $ \xi_1, \xi_2, \dotsc $, а
    $ g_n(t) $ "--- характеристическая функция надлежащим образом нормированной суммы
    $ \frac{S_n}{\sqrt n} $ этих величин:
    $$ g_n(t) = \Bigl( f \bigl( \frac{t}{\sqrt n} \bigr) \Bigr)^n $$
    Учитывая равенства,
    $$ f'(0) = \ii \mathbb E \xi = 0, \quad f''(0) = -\mathbb \xi^2 = -1, $$
    в окрестности нуля справедливо
    $$ f(t) = f(0) + tf'(0) + \frac12 t^2 f''(0) + o(t^2) = 1 + \ii t \mathbb E \xi - \frac{t^2}2
    \mathbb E \xi^2 + o(t^2) = 1 - \frac{t^2}2 + o(t^2) $$

    Рассмотрим при $ n \to \infty $, асимптотику функции
    $$ \ln g_n(t) = b \ln f \Bigl( \frac t {\sqrt n} \Bigr) $$
    При фиксированных $ t $ можно написать, что
    $$ f \Bigl( \frac{t}{\sqrt n} \Bigr) = 1 - \frac{t^2}{2n} + o \Bigl( \frac1n \Bigr), \quad
    n \to \infty $$
    При этом, при $ \alpha \to 0 $ справедливо
    $$ \ln(1 + \alpha) = \alpha - \frac{\alpha^2}2 + o(\alpha^2) $$
    Значит,
    $$ \ln f \Bigl( \frac{t}{\sqrt n} \Bigr) = -\frac{t^2}{2n} + o \Bigl( \frac1n \Bigr), \quad
    n \to \infty $$
    $$ \ln g_n(t) = -\frac{t^2}{2n} + o \Bigl( \frac1n \Bigr), \quad n \to \infty $$
    $$ g_n(t) \underarr{n \to \infty} e^{-\frac{t^2}2} $$
    Справа получили характеристическую функцию нормального распределения $ N(0, 1) $.
\end{proof}

Приведём также один из более общих вариантов ЦПТ:
\begin{theorem}[Ляпунова]
    $ \xi_1, \xi_2, \dotsc $ "--- независимые случайные величины, $ \quad
    S_n = \xi_1 + \dots + \xi_n $ "--- их суммы, $ \quad a_k = \mathbb E \xi_k, \quad
    \sigma_k^2 $ "--- их дисперсии, $ \quad \gamma_{3k} = \mathbb E |\xi_k -
    \mathbb E \xi_k|^3 $ "--- центральные моменты третьего порядка
    $$ A_n = a_1 + \dots + a_n = \mathbb E S_n, \quad
    B_n^2 = \mathcal D S_n = \sigma_1^2 + \dots + \sigma_n^2, \quad
    \Gamma_n = \sum_{k = 1}^n \gamma_{3k} $$
    Обозначим \emph{дробь Ляпунова}
    $$ L_n = \frac{\Gamma_n}{B_n^3} $$
    $$ \implies \sup\limits_x \Bigl| P\Set{\frac{S_n - A_n}{B_n} < x} - \Phi(x) \Bigr| \le CL_n, $$
    где $ C $ "--- абсолютная (не зависящая от случайных величин и $ n $) константа.
\end{theorem}

Важную роль в теории вероятностей имеют функция распределения и плотность распределения
стандартного нормального распределения $ N(0, 1) $:
$$ \Phi(x) = \frac1{\sqrt{2\pi}} \int_{-\infty}^x e^{-\frac{t^2}2}\di t =
\int_{-\infty}^x \phi(t) \di t, \quad \phi(x) = \frac1{\sqrt{2\pi}} e^{-\frac{x^2}2} $$

\section{Законы больших чисел}

\begin{definition}
    $ \xi_1, \xi_2, \dotsc $ "--- последовательность случайных величин с математическими
    ожиданиями $ a_1, a_2, \dotsc $

    Говорим, что эта последовательность удовлетворяет \emph{закону больших чисел}, если
    $$ \frac{\xi_1 + \dots + \xi_n}n \xrightarrow[n \to \infty]{P} \frac{a_1 + \dots + a_n}n $$
\end{definition}

\begin{theorem}[Чебышёва]
    $ \xi_1, \xi_2, \dotsc $ "--- независимые случайные величины с математическими ожиданиями
    $ a_1, a_2, \dotsc $ и дисперсиями $ \sigma_1^2, \sigma_2^2, \dotsc $.
    Если $ \sigma_k^2 \le c $ для некоторого $ c $, то рассматриваемая последовательность
    \textbf{удовлетворяет закону больших чисел}.
\end{theorem}

Теорема Чебышёва представляет собой частный случай следующего утверждения:
\begin{theorem}[Маркова]
    Для произвольных случайных величин $ \xi_1, \xi_2, \dotsc $ выполняется
    $$ \frac{\mathcal D(\xi_1 + \dots + \xi_n)}{n^2} \underarr{n \to \infty} 0 $$

    Тогда эта последовательность \textbf{удовлетворяет закону больших чисел}.
\end{theorem}

\begin{theorem}[Хинчина]
    $ \xi_1, \xi_2, \dotsc $ "--- независимые случайные величины с одинаковым распределением
    и конечным математическим ожиданием $ a $.

    $$ \frac{\xi_1 + \dots + \xi_n}n \xrightarrow[n \to \infty]{P} a $$
\end{theorem}

\begin{remark}
    Из теоремы Хинчина следует соотношение для числа успехов $ \mu_n $ в схеме Бернулли при
    вероятности успеха $ p $ в каждом испытании:
    $$ \frac{\mu_n}n \overarr{P} p $$
\end{remark}

\section{Закон нуля и единицы и усиленный закон больших чисел}

Пусть $(\Omega, F, P) $ "--- вероятностное пространство, $ A_1, A_2, \dotsc $ "--- произвольная
последовательность событий из $ \sigma $-алгебры $ F $.
Рассмотрим событие
$$ A = \bigcap_{n = 1}^\infty \bigcup_{k = n}^\infty A_k $$
$ A $ означает, что появляется бесконечное число событий в данной последовательности.

\begin{lemma}[Бореля"--~Кантелли]
    \hfill
    \begin{enumerate}
        \item Если $ \sum P\Set{A_n} < \infty $, то $ P\Set{A} = 0 $.
        \item Если события $ A_1, A_2, \dotsc $ независимы и $ \sum P\Set{A_n} = \infty $, то
            $ P\Set{A} = 1 $.
    \end{enumerate}
\end{lemma}

\begin{theorem}[закон нуля и единицы]
    $ A_1, A_2, \dotsc $ "--- независимые события
    $$ P\Set{A} =
    \begin{cases}
        0, \quad \sum P\Set{A_n} < \infty \\
        1, \quad \sum P\Set{A_n} = \infty
    \end{cases} $$
\end{theorem}

\begin{definition}
    Последовательность случайных величин $ \xi_1, \xi_2, \dotsc $ удовлетворяет \emph{усиленному
    закону больших чисел}, если существует последовательность чисел $ \beta_n $, для которой верно
    $$ \frac{\xi_1 + \dots + \xi_n}n \xrightarrow[n \to \infty]{\text{п. н.}} \beta_n $$
\end{definition}

\begin{theorem}[Колмогорова]
    $ \xi_1, \xi_2, \dotsc $ "--- последовательность независимых одинаково распределённых случайных
    величин.

    Соотношение
    $$ \frac{\xi_1 + \dots + \xi_n}n \overarr{\text{п. н.}} a $$
    выполняется \textbf{тогда и только тогда}, когда у исходных случайных величин существует
    конечное математическое ожидание, к которому нормированные суммы и стремятся.
\end{theorem}

\begin{theorem}[Колмогорова]
    $ \xi_1, \xi_2, \dotsc $ "--- независимые случайные величины с математическими ожиданиями $ a_1,
    a_2, \dotsc $ и дисперсиями $ \sigma_1^2, \sigma_2^2, \dotsc $
    $$ \sum_{k = 1}^\infty \frac{\sigma_k^2}{k^2} < \infty $$

    Тогда последовательность $ \xi_1, \xi_2, \dotsc $ удовлетворяет УЗБЧ.
\end{theorem}
